\documentclass[11pt,a4paper]{article}
\usepackage[T1]{fontenc}
\usepackage{lmodern,microtype,amsmath,amssymb,amsthm,booktabs,array}
\usepackage[a4paper,margin=26mm,top=23mm,bottom=25mm,headheight=14pt]{geometry}
\usepackage[dvipsnames]{xcolor}
\usepackage{fancyhdr}
\usepackage[colorlinks=true,allcolors=MidnightBlue,pdfauthor={Keio research notes},pdftitle={Divisibility and residue patterns in the Gilbreath prime triangle}]{hyperref}
\definecolor{inkblue}{HTML}{194C64}
\newtheorem{theorem}{Theorem}[section]
\newtheorem{proposition}[theorem]{Proposition}
\newtheorem{corollary}[theorem]{Corollary}
\pagestyle{fancy}\fancyhf{}\fancyhead[L]{\small\sffamily Gilbreath's prime triangle}
\fancyhead[R]{\small\sffamily November 2026}\fancyfoot[C]{\thepage}
\renewcommand{\headrulewidth}{0.3pt}
\setlength{\emergencystretch}{2em}
\raggedbottom
\begin{document}
\thispagestyle{plain}
\noindent{\small\sffamily\bfseries\color{inkblue} NUMBER THEORY \quad / \quad RESEARCH NOTE}\par
\vspace{0.65em}
\noindent{\LARGE\bfseries Divisibility and residue patterns in the\\[2pt] Gilbreath prime triangle}\par
\vspace{0.5em}
\noindent{\normalsize Western tails, interior rays, and the realization of finite extensions}\par
\vspace{0.35em}
\noindent{\small 21 November 2026}\par
\vspace{0.7em}
\begin{abstract}
For the triangle of iterated absolute differences of consecutive primes, the greatest common divisor of every tail of the western edge is one. On every fixed interior column, even after sampling depths at multiples of any fixed power of two, the two residues zero and two modulo four both recur infinitely often. The associated binary series is transcendental over the rational function field of characteristic two. Along the moving right edge of the first $n$ primes, the number of entries in each of those residue classes tends to infinity. For the number $N_n$ of length-$n$ Gilbreath prefixes we prove $\log N_n=\Theta(n^2)$. We also give a binomial parity criterion and a constructive solution to the small-width right-edge realization question for finite Gilbreath prefixes. Finally, a separate constant-edge conjecture of Kimberling for primes in arithmetic progressions admits a counterexample at modulus seven and fails at the first difference for a density-one set of odd moduli. The classical assertion that every positive-depth western entry equals one remains open.
\end{abstract}
\section{The prime conjecture and the available information}
Write $p_j$ for the $j$th prime, beginning with $p_1=2$. For $j\ge1$ and $r\ge0$ put
\[
 A_{0,j}=p_j,\qquad A_{r+1,j}=|A_{r,j+1}-A_{r,j}|.
\]
Proth's 1878 question \cite{proth}, usually called Gilbreath's conjecture, asks whether
\begin{equation}\label{eq:gil}
                         A_{r,1}=1\qquad\text{for all }r\ge1 .
\end{equation}
The first gap is $1$. The rest of the first difference row, starting at $j=2$, consists of even numbers. Two directions in the triangle will matter. In column $j$, the depth $r$ varies with $j$ fixed. In the right edge of the prefix ending at $p_n$, the entries are $A_{r,n-r}$ with $n$ fixed. For example, for the prefix $(2,3,5,7,11,13)$ the four interior entries of its right edge, in order of increasing depth, are $2,2,0,2$.

Odlyzko verified the conjecture through the primes below $10^{13}$ \cite{odlyzko}. Colonna reports a calculation through $1.5\times10^{15}$ \cite{colonna}. Both bounds specify prime magnitude. Chase, Hunter and Tao prove a Cram\'er-model analogue and, assuming a Cram\'er-type prime-gap bound, give a deterministic inverse theorem involving long blocks of zeroes or long shallow two-valued blocks \cite{cht}. Neither result supplies the universal bound on column two needed in \eqref{eq:gil}. Bhat, Cobeli and Zaharescu's Conjecture~1 asks for limiting frequencies of $1/2$ on the interior rays modulo four \cite{bcz}. Our recurrence assertions give no limiting frequencies.

We searched the 225 numbered TeX manuscripts at the repository revision recorded in \cite{vmthree}. The directly relevant sections occur in manuscripts 003, 048 and 072. The last two TeX files are byte-identical at this revision. Manuscript 003 uses the parity identity for iterated differences to analyze bounded-height obstruction words. Manuscript 048 uses the right edge and reverse absolute differences to classify small widths, and leaves the realization of its candidate edges at prescribed lengths open \cite{vmwidth}. Here the parity calculation incorporates long strings of consecutive primes in prescribed residue classes. The localized argument in Section~\ref{ray:sec} uses Maynard's theorem. The finite construction in Section~\ref{sec:reach} uses the fact that on binary data a parity calculation recovers the entire integer triangle. These two mechanisms have different consequences for the original question.

\section{Divisibility on the western edge}\label{sec:west}
Throughout this section $p_1=2$, $A_{0,j}=p_j$, and
$A_{r+1,j}=|A_{r,j+1}-A_{r,j}|$. Write $w_r=A_{r,1}$.
For every $r\ge1$, $w_r$ is odd and $A_{r,j}$ is even for $j\ge2$.
This follows first for the gaps and then by induction.
Since $w_1=1$ and $w_{r+1}=|w_r-A_{r,2}|$, the original assertion
$w_r=1$ for every $r\ge1$ is equivalent to
\[
                         A_{r,2}\in\{0,2\}\quad(r\ge1).
\]
Indeed, if $w_r=w_{r+1}=1$, the only possibilities for $A_{r,2}$ are
$0$ and $2$. The converse follows by induction. If $w_s$ is the first
failure, then $A_{s-1,2}\ge4$ and $w_s=A_{s-1,2}-1$.

\begin{theorem}[Gcd of each vertical tail]
For all $R,j\ge1$,
\[
 \gcd\{A_{r,j}:r\ge R\}=
 \begin{cases}1,&j=1,\\2,&j\ge2.\end{cases}
\]
In particular, the only possible eventual constant on the western edge is $1$.
\end{theorem}
\begin{proof}
Suppose $q\ge2$ divides $A_{r,j}$ for every $r\ge R$.
Whenever $q\mid u$ and $q\mid |u-v|$, also $q\mid v$.
Apply this observation to every adjacent pair of heights in column $j$
and then repeat it in each successive column. It follows that
$q\mid A_{R,i}$ for every $i\ge j$.

For each $s\ge1$ there is an $i\ge j$ with $A_{s,i}\equiv0\pmod q$.
For $q\ge3$, Shiu's theorem supplies $s+1$ consecutive primes in
the same reduced residue class modulo $q$, arbitrarily far out
\cite{shiu}. All gaps in that block, and then the entire triangle
below those gaps, are divisible by $q$. For $q=2$ the assertion follows
from any block of odd primes. If now $s\ge2$ and $q$ divides all
$A_{s,i}$ for $i\ge j$, their predecessors $A_{s-1,i}$, $i\ge j$,
are constant modulo $q$. The zero just obtained makes that common
residue zero. Descending to $s=1$, we find that $q$ divides every
prime gap $p_{i+1}-p_i$ for $i\ge j$.

At $j=1$ the first gap is $1$. At $j\ge2$, divisibility of all the
gaps forces all primes past $p_j$ to occupy one residue class modulo
$q$. If $q>2$, Dirichlet's theorem gives infinitely many primes in
each of the two distinct reduced classes $1,-1\pmod q$, a
contradiction. Conversely, all entries in every column $j\ge2$ at
positive height are even. These observations also exclude the
possibility of an all-zero tail.
\end{proof}

The argument gives a finite witness for failure of any proposed common
divisor of a western tail. Fix $R\ge1$ and $q\ge2$, and take any $m\ge2$
for which $p_m,\ldots,p_{m+R-1}$ all agree modulo $q$.
Then
\[
    q\nmid w_t\quad\hbox{for at least one }t\in[R,R+m-1].
\]
Otherwise divisibility on this interval propagates to
$A_{R,1},\ldots,A_{R,m}$. Moreover $q\mid A_{s,m}$ for
$1\le s<R$ by the chosen prime block. Descending one height at a
time, congruence of the entries in positions $1,\ldots,m$ together
with this anchor at $m$ yields divisibility throughout those positions
at every positive height down to $1$. This contradicts $A_{1,1}=1$.
The assertion at $R=1$ follows directly from that same gap.

\paragraph{What the modular information leaves open.}
Consider the strictly increasing sequence
\[
           a_1=2,\qquad a_2=3,\qquad a_i=p_i+4\quad(i\ge3).
\]
Its initial gaps are $1,6$, so its western value at height $2$ is $5$.
Nevertheless all its odd entries have the same residues modulo $4$
as the corresponding primes. For any $j\ge2$ the entire interior ray
therefore has exactly the prime ray's residues modulo $4$ at all positive
heights. To see this, subtract $1$ from the odd top entries, divide
by $2$, and reduce modulo $2$, where $|u-v|\equiv u+v$.
Every sufficiently distant block of congruent primes remains a block of
congruent integers after the common translation by $4$.
The same descent proof, using the unchanged first gap $1$, gives gcd
$1$ for every tail of this sequence's western edge. Thus even complete
knowledge of all the interior residues modulo $4$ together with the
western tail gcds cannot recover the magnitude $A_{r,2}\le2$ without
further arithmetic information about the actual primes.

The example can preserve an arbitrarily long prescribed prime prefix.
Fix $L\ge2$ and leave $p_1,\ldots,p_L$ unchanged, translating all
subsequent primes by the same positive multiple $C$ of $4$. Put
$e_0=p_L$ and $e_s=A_{s,L-s}$ for $1\le s<L$.
For $p_{L+1}+C>2+\sum_{s=0}^{L-1}e_s$, each new entry on the
right boundary of the enlarged finite triangle is obtained by
subtracting the corresponding $e_s$ from the preceding new entry.
Consequently its western value at height $L$ is precisely
$p_{L+1}+C-\sum_{s=0}^{L-1}e_s>2$.
All preceding heights agree with the original triangle, all interior
residues modulo $4$ still agree at every height, and the argument
for the tail gcd continues to apply.

One can require all terms of the comparison sequence to be prime.
For $i>L$ choose successively an arbitrarily large prime
$a_i\equiv p_i\pmod{4i!}$ with $a_i>a_{i-1}$, taking
$a_{L+1}>2+\sum_{s=0}^{L-1}e_s$. The choices exist by Dirichlet's
theorem since $p_i>i$ and $p_i$ is odd. For every fixed modulus
$q$ one has $a_i\equiv p_i\pmod q$ for all $i\ge\max(q,L+1)$.
Hence the required congruent blocks exist for every $q$, all interior
rays agree modulo $4$, and every western tail has gcd $1$. The western
value at height $L$ again exceeds $2$. This construction omits intervening
primes, which is precisely where the comparison departs from the original
problem.

\paragraph{A sharp threshold for the infinitely-many-ones question.}
If $A_{r,2}\le6$ at all sufficiently large heights of the prime triangle,
then $w_r=1$ infinitely often. Otherwise $w_r\ge3$ eventually, so that
$0\le A_{r,2}\le2w_r$ and $w_{r+1}=|w_r-A_{r,2}|\le w_r$ eventually.
This nonincreasing sequence of integers stabilizes, and its constant must
be $1$ by the tail gcd theorem. Equivalently, finitely many western ones
would require $A_{r,2}\ge8$ infinitely often. The bound $6$ is sharp
for the recurrence together with the western tail gcd condition. The
strictly increasing sequence $2,3,15,19,23,27,\ldots$ begins with the
gap row $(1,12,4,4,\ldots)$ and the next row $(11,8,0,0,\ldots)$.
Thereafter the second entry stays $8$, and the western sequence of
positive-height entries is $1,11,3,5,3,5,\ldots$. Its every tail has
gcd $1$.

\section{Transcendental series on interior rays of the prime triangle}
\label{ray:sec}

Write $p_j$ for the $j$th prime, with $p_1=2$, and set
\[
 A_{0,j}=p_j,\qquad A_{k+1,j}=|A_{k,j+1}-A_{k,j}|.
\]
For $j\ge2$, every $A_{k,j}$ with $k\ge1$ is even. Bhat, Cobeli and Zaharescu study these fixed-column rays and conjecture a quantitative balance of the \emph{exact} values $0$ and $2$ \cite[Conjecture~2]{bcz2}. We obtain an unconditional arithmetic property of each ray, including samples at every fixed power-of-two spacing. For $h\ge0$ set $d=2^h$ and define integers
\[
 a_0=(p_j-1)/2,\qquad a_m=A_{dm,j}/2\quad(m\ge1).
\]
The dependence of $a_m$ on $j,d$ is understood throughout this section.

\begin{theorem}[Interior rays at dyadic heights]\label{ray:main}
For every $j\ge2$ and every $d=2^h$, the series $\sum_{m\ge0}(a_m\bmod 2)X^m$ is transcendental over $\mathbb F_2(X)$. In particular infinitely many $A_{dm,j}$ belong to each of the residue classes $0$ and $2$ modulo $4$. Both complex series
\[
 \sum_{m\ge1}(a_m\bmod 2)z^m
 \qquad\hbox{and}\qquad
 \sum_{m\ge1}A_{dm,j}z^m
\]
have the unit circle as a natural boundary. Here a bit in a complex series denotes its representative in $\{0,1\}$.\end{theorem}

\paragraph{A local consequence of Maynard's theorem.}
For each fixed $s\ge1$ and $a\in\{1,3\}$, there is a constant $C_s$ such that for every sufficiently large $Y$ there are $s$ consecutive primes congruent to $a\pmod4$ in $[Y,Y+C_s Y^{2/3}]$. We give the localization argument because strings known to occur infinitely often \cite{shiu,bftb} alone do not control their absence from a specified interval.

By Maynard's short-interval theorem \cite[Theorem~3.4]{maynard}, for a fixed $s$ there exists $k$ such that any fixed admissible $k$-tuple of distinct positive linear forms has at least $s$ prime values at the same integer in $[x,x+x^{2/3}]$ for every sufficiently large $x$. Indeed, take the fixed exponent parameter $1/24$ in that theorem, so $7/12+1/24<2/3$. Put $h_i=4k!i$ for $1\le i\le k$. For each integer $t\in[h_1,h_k]$ other than the $h_i$, choose a distinct prime $r_t>h_k$. With $M=4\prod_t r_t$, the Chinese remainder theorem gives a positive $b$ with $b\equiv a\pmod4$ and $b\equiv-t\pmod{r_t}$ for every such $t$. The forms $Mn+b+h_i$ are admissible. A prime dividing $M$ divides none of the $b+h_i$. At a prime $\ell\nmid M$ with $\ell\le k$ the $h_i$ coincide modulo $\ell$ and forbid just one residue of $n$, and at $\ell>k$ at most $k$ residues are forbidden. For large $n$, each $Mn+b+t$ at an unselected offset is composite. Consequently the prime values at the selected offsets, taken in order, are consecutive primes and all are $a\pmod4$. Apply the short-interval theorem with $x=\lceil Y/M\rceil$ and enlarge $C_s$ to include the fixed offsets and scaling factor. This proves the assertion.

\paragraph{A nonautomatic word on each progression of prime indices.}
For every fixed integer $D\ge1$, put
\[
 u_n=\big((p_{j+Dn}-1)/2\big)\bmod 2,\qquad n\ge0.
\]
This word is not automatic in any integer base $b\ge2$. To prove it, fix a bit $e$. Shiu's theorem supplies arbitrarily long strings of consecutive primes in the residue class giving $e$ \cite{shiu}. Intersecting such a string of length $DL$ with the index progression produces a block $e^L$. In the other direction, for any fixed $C>0$ and all sufficiently large $L$, no block $e^L$ can start at an index $i\le CL$. If it did, put $P=p_{j+Di}$ and $Q=p_{j+D(i+L-1)}$. The prime number theorem gives $Q=O_{C,D,j}(L\log L)$, while $Q-P\ge D(L-1)$. Apply the local assertion with $s=D$ to the opposite residue class starting at $Y=\max(P+1,Y_0)$ for a fixed sufficiently large threshold $Y_0$. Then $Y+C_DY^{2/3}<Q$ for all sufficiently large $L$, uniformly in $i\le CL$. Among these $D$ consecutive primes is one with index congruent to $j$ modulo $D$, contradicting $e^L$.

The automatic run obstruction also follows from Schlage-Puchta's criterion \cite[Theorem~3]{puchta}. For completeness, a $b$-automatic word has a constant $C$ such that every factor of length $L$ appears at an index $\le CL$. In the standard finite-state presentation it is a letterwise coding of a fixed point of a $b$-uniform morphism $\sigma$. This presentation follows by allowing leading zeroes at the start state and replacing each state by the list of its $b$ digit-successors. Choose $b^{v-1}<L\le b^v$. Every factor in the coded word lies in the coding of $\sigma^v(AB)$ for an occurring pair of adjacent states $AB$. Since there are finitely many such pairs, all have an occurrence before a fixed index $K$. The factor therefore occurs before $(K+2)b^v\le b(K+2)L$. This proves the claim and the nonautomaticity of $u$. For $D=1$ the same local exclusion follows directly from the prime number theorem in the two progressions modulo $4$, which gives $N/2+o(N)$ ones among the first $N$ terms. The number of ones between indices $i$ and $i+L-1$ is then $L/2+o(L)$ uniformly for $i\le CL$.

\paragraph{Passing to the ray.}
Work over $\mathbb F_2$, and let $E$ be the forward shift of the prime index. Reduction modulo $2$ changes an absolute difference into addition. Applying this to $(p_i-1)/2$ for $i\ge2$, whose positive absolute differences are exactly half those in the prime triangle, gives
\[
 a_m\bmod2=\sum_{t=0}^m\binom mt u_t\pmod2,\qquad D=d.
\]
Indeed $(1+E)^{dm}=(1+E^d)^m$. Write $U(X)=\sum u_t X^t$ and $V(X)=\sum (a_m\bmod2)X^m$. Summing the displayed identity yields
\[
 V(X)=\frac{1}{1+X}U\!\left(\frac{X}{1+X}\right).
\]
The transformation on the right is an involution and carries algebraic formal power series to algebraic formal power series. By Christol's theorem \cite{christol}, algebraicity of $U$ over $\mathbb F_2(X)$ is equivalent to $2$-automaticity of $u$. The preceding argument therefore gives transcendence of $U$ and $V$. A binary sequence assuming just one of its two values from some point on has a rational series. This proves the two infinitude assertions.

\paragraph{Passing to complex coefficients.}
Both the bit sequence and the integer sequence $(a_m)$ have integer coefficients and reduce to $V$ modulo $2$. Neither of their complex series can be rational. Indeed, a rational series has finite Hankel rank. Its integer Hankel minors above that rank vanish and hence also vanish after reduction, giving a rational series over $\mathbb F_2$. The finite-rank criterion follows by taking a linear dependence of initial Hankel columns, which is a constant-coefficient recurrence. Moreover, $0\le A_{k,j}\le p_{j+k}$, since each absolute difference is at most the larger parent. Thus $a_m=O_{d,j}(m\log(m+2))$, and infinitely many $a_m$ are odd. The two complex series have radius of convergence exactly one. The P\'olya--Carlson theorem \cite{pc} now gives their natural boundaries. Multiplication by $2$ and a change of constant coefficient recover the asserted integer series.

This conclusion concerns residues on interior rays. For instance, an entry equal to $6$ contributes to the residue class $2\pmod4$. Neither the exact $0,2$ discrepancy on an interior ray nor the value of every entry of the western edge follows from the series theorem.

\begin{corollary}[Finite linear combinations of columns]
Fix $d=2^h$, $j\ge2$, and rational $c_0,\ldots,c_t$ not all zero. Then
\[
 \sum_{m\ge1}\left(\sum_{s=0}^t c_s A_{dm,j+sd}\right)z^m
\]
has the unit circle as a natural boundary.
\end{corollary}
\begin{proof}
Scale to integral coefficients with at least one odd. In the proof above the source word for the combination of the half entries is $\sum_s(c_s\bmod2)u_{m+s}$. Its binary series equals $\big(\sum_s(c_s\bmod2)X^{-s}\big)U(X)$ plus a Laurent polynomial. The multiplier is nonzero, so this series remains transcendental over $\mathbb F_2(X)$. The Pascal involution preserves transcendence. Infinitely many coefficients of the half-normalized integer combination are thus odd, and their absolute magnitudes are $O(m\log(m+2))$. The same Hankel-rank and Pólya--Carlson arguments apply.
\end{proof}

\section{Variation on the right edge of every long prime prefix}

Write $p_1=2,p_2=3,\ldots$ for the primes in increasing order and
$A_{r,j}=(D^rp)_j$, where $(Dx)_j=|x_{j+1}-x_j|$.
For $n\ge3$ let
\[
 R(n)=\#\{1\le r\le n-2:A_{r,n-r}\equiv2\pmod4\}.
\]
Every entry counted here belongs to the original prime array.
The range ends one cell before its left edge.

\begin{theorem}[Valuation bands] For every integer $h\ge1$ there exist
$0=t_0<t_1<\cdots<t_h$ and $N_h$ such that for each $n\ge N_h$
and each $0\le j<h$ there is an $r$ counted by $R(n)$ with
$t_j\le v_2(r)<t_{j+1}$. In particular $R(n)\to\infty$.\end{theorem}

\begin{proof}
For $0\le i\le n-2$ let
$f(i)=(p_{n-i}-1)/2\pmod2$ and expand in $\mathbf F_2$ as
\[
 f(i)=\sum_{r=0}^i \binom{i}{r}a_r,\qquad
 a_r=\sum_{j=0}^r \binom{r}{j} f(j).
\]
The inversion follows from
$\sum_j\binom{r}{j}\binom{j}{k}=\binom{r}{k}2^{r-k}$.
For $r\ge1$ the whole difference cone uses odd primes. Subtracting 1
from each input and dividing by 2 divides its $r$th absolute difference
by 2. Since absolute and ordinary differences agree modulo 2, and
reversal preserves the apex of a triangle of length $r+1$, we have
$a_r\equiv A_{r,n-r}/2\pmod2$. The binary factorization of
$(1+X)^i$ shows that $\binom{i}{r}$ is odd precisely when every 1-bit of
$r$ is a 1-bit of $i$.

Fix $t\ge0$. By Shiu's theorem \cite{shiu}, one can choose two disjoint
finite intervals of indices among the odd primes, each containing a run
of at least $2^{t+1}$ consecutive primes congruent to 1 modulo 4 and
a run of that length congruent to 3 modulo 4. Choose $T>t$ such that
the interval spanning both selected intervals has length less than $2^T$.
Once $p_n$ includes this entire prefix, reflect indices via $j\mapsto n-j$.
At most one of the two reflected intervals can straddle a boundary
between consecutive aligned blocks of $2^T$ indices. In the other interval,
choose multiples $i,i'$ of $2^t$, one in each run. Their indices have
identical bits in positions $\ge T$, all bits below $t$ vanish, and
$f(i)\ne f(i')$. At least one term with $a_r=1$ must distinguish these
indices. Its bits below $t$ vanish, while it has some bit in positions
$t,\ldots,T-1$. Thus $t\le v_2(r)<T$ and $r\ge1$.
Starting with $t_0=0$, repeat this step $h$ times and take the maximum
of the finitely many prefix thresholds. The bands are disjoint.
\end{proof}

The same argument applies to any binary input having arbitrarily long runs
of each symbol in every tail. Shiu supplies precisely this input for the two
odd residue classes modulo 4. The proof gives no bound on the leftmost
entry $A_{n-1,1}$.

\begin{corollary}[Distance from periodic words]\label{diag:mask}
For the same bands, there is a constant $J_h$ such that the following
holds uniformly for every power of two $P$ and every $n\ge J_h+P$.
Given any binary sequence $(w_r)_{r\ge0}$ of period $P$, each band
contains a depth $1\le r\le n-2$ with
$A_{r,n-r}/2\not\equiv w_r\pmod2$.
In particular, each band contains both a depth with
$A_{r,n-r}\equiv0\pmod4$ and one with $A_{r,n-r}\equiv2\pmod4$
once $n\ge J_h+1$.\end{corollary}

\begin{proof}
Take $J_h$ to be the largest prime index used in any of the $h$ pairs
of witness intervals above. Over $\mathbf F_2$, define
$d_i=\sum_{r=0}^i\binom{i}{r}w_r$. With $E$ the forward shift,
$d_i=((1+E)^iw)_0$. Since $P$ is a power of two,
$(1+E)^Pw=(1+E^P)w=0$, so $d_i=0$ for $i\ge P$.
All witness indices $i=n-j$ lie in this range when $n\ge J_h+P$.
Thus $f+d$ agrees with $f$ on every witness interval. Pascal inversion
gives $a_r+w_r$ as its coefficients. The proof of the theorem supplies a
coefficient equal to 1 in each band. Taking $w$ successively constant
zero and constant one gives the last assertion.
\end{proof}

\paragraph{A uniform numerical lower bound.}
There is a short finite certificate for the statement that, for every $n\ge4097$, at least $952$ entries in the range $1\le r\le n-2$ are $0\bmod4$ and at least $952$ are $2\bmod4$. The bound applies to arbitrarily large prime prefixes.

Put $x_j=(p_j-1)/2\bmod2$ for $j\ge2$. For a power of two $B$, take the first $2B-1$ symbols $x_2,\ldots,x_{2B}$. For each of their $B$ consecutive length-$B$ windows, reverse the word and compute its binary Pascal transform. Let $W_B$ be the minimum number of ones in these transforms.

For $n\ge2B+1$, the corresponding $2B-1$ indices in $f$ form an interval contained in $[1,n-2]$. It contains one aligned block $[vB,(v+1)B-1]$. On this block each monomial $\binom{i}{r}\bmod2$ restricts to zero or to a single monomial in the $\log_2 B$ low bits of $i$. Restriction therefore cannot increase the number of nonzero coefficients. The transform of either constant mask $w$ vanishes at all $i\ge1$, as in Corollary~\ref{diag:mask}. Applying the restriction to $f+d$ thus gives at least $W_B$ global coefficients where $a_r\ne w_r$. Discarding $r=0$ costs at most one. Each of the two residue counts is at least $W_B-1$.

\[
\begin{array}{c|rrrrrrrr}
B&16&32&64&128&256&512&1024&2048\\\hline
W_B&6&12&25&48&103&221&454&953
\end{array}
\]
The following Python function reproduces each integer using trial division and the binary subset transform. It needs no external packages.
\begin{verbatim}
from math import isqrt

def certificate(B):
    primes, x = [], 2
    while len(primes) < 2*B:
        if all(x % j for j in range(2, isqrt(x)+1)):
            primes.append(x)
        x += 1
    bits = [((p-1)//2) % 2 for p in primes[1:]]
    def weight(offset):
        v = bits[offset:offset+B][::-1]
        stride = 1
        while stride < B:
            for i in range(B):
                if i & stride:
                    v[i] ^= v[i ^ stride]
            stride *= 2
        return sum(v)
    return min(weight(i) for i in range(B))
\end{verbatim}

\section{Which small-width right edges occur?}
\label{sec:reach}
For a strictly increasing Gilbreath prefix $S=(s_1,\ldots,s_n)$ with $(s_1,s_2)=(2,3)$, let $K_S$ be the set of all integer extensions preserving the left edge. An extension in $K_S$ may be smaller than $s_n$. As in Muney~\cite{muney}, write $e_i$ for the last entry of row $i$, $1\le i<n$, and put $e^0=(2,0,\ldots,0,1)$, with $n-1$ entries. We assume $n\ge3$. Corpus manuscript \cite{vmwidth}, Part B, subsection on widths at most nine, proves that
\[
 |K_S|\le9\quad\Longleftrightarrow\quad
 e=e^0+2\mathbf 1_P,\qquad P\subseteq\{1,\ldots,n-2\},\quad |P|\le2,
\]
where the indicator is padded by a final zero. In this case $|K_S|=5+2|P|$. Its subsequent remark asks which such $P$ can actually occur. This is also the realization step in the stability problem posed by Muney in Section~16(5).


\paragraph{Why the width cutoff is exact.}
Here is the argument from \cite{muney,vmwidth} in the notation of this section. In any positive row, the entry after its initial $1$ is $0$ or $2$, since its difference from $1$ is itself the initial $1$ below. From $A_{i,j+1}\equiv A_{i,j}+A_{i+1,j}\pmod2$, induction across the columns shows that every entry after the first in a positive row is even. In particular the first gap is $1$ and all later gaps are even, so all top entries after $2$ are odd.

For a proposed extension, start with $d=|k-s_n|$ and successively apply $x\mapsto |x-e_i|$ for $i=1,\ldots,n-1$. Only the new apex can violate the left-edge condition, so its value must be $1$. Reading this recursion backwards from $\{1\}$ first gives $\{0,2\}$, since $e_{n-1}=1$. An even fold of size $e$ sends a distance set $T$ to
\[
     (e+T)\ \cup\ \{e-u:u\in T,\ u\le e\}.
\]
If $T=[0,M]\cap2\mathbb Z$ and $e\le M$, the result is another full even interval, with endpoint $M+e$. Before the first positive even fold, all intervening rows end in zero. Reversing the rows one at a time shows these rows are $(1,0,\ldots,0)$, and the first row ending in a positive even number has the form $(1,2,\ldots,2)$. Indeed the right endpoint specifies the constant interior forced by a row of zero differences; the first difference fixes its value to $0$ or $2$. Thus the first positive even fold is $2$. At a first fold with $e>M$ one must already have $M\ge4$; the resulting set has $M+1\ge5$ elements and omits zero. Later folds never reduce its size. While a set omits zero, the first fold that creates zero adds it to the translated copy as a new element. Consequently the final two-sided signed distance set has at least ten elements if a fold ever exceeds the current endpoint.

Writing $A=\sum_i e_i$, when no fold exceeds the endpoint the two-sided extension width is $A+2$. It is always at most $A+2$. It follows that width at most nine is equivalent to $A\le7$. The vector $e-e^0$ has nonnegative even entries of total at most four. If one entry were four, all the others would vanish. Below that entry, all even endpoints would be zero, forcing its endpoint to be at most two by the row argument above. This excludes both an entry four at a later level and an excess four in the first gap. The asserted list and widths follow.

\begin{theorem}[Realizable right edges]
\label{thm:reachable}
Let $m=n-2$. For any $P\subseteq\{1,\ldots,m\}$, set $b_k=\mathbf1_P(k+1)$ for $0\le k<m$. The integer edge $e^0+2\mathbf1_P$ occurs if and only if
\[
 \sum_{k=0}^{m-1}\binom{m-1}{k}b_k\equiv0\pmod2.
\]
When this holds, there is a unique realization in which every gap after $3$ is either $2$ or $4$. Explicitly define its entries by $s_1=2$, $s_2=3$, $s_{i+2}=s_{i+1}+2(1+c_i)$, where $c_{m-j}$ is the representative in $\{0,1\}$ of
\[
 \sum_{k=0}^j\binom{j}{k}b_k\pmod2,\qquad 0\le j<m.
\]
\end{theorem}
\begin{proof}
Let $\Delta$ be unsigned adjacent differencing of a finite word. All entries $s_i$, $i\ge2$, are odd for a Gilbreath prefix by the parity propagation in \cite{muney}. The row-two condition gives $s_3=5$ since the positive even gap after $3$ differs from $1$ by $1$. Thus the normalized excess gaps $c_i=(s_{i+2}-s_{i+1})/2-1$ are nonnegative integers and $c_1=0$. Their integer right-edge trace is $b_k=(\Delta^k c)_{m-k}$ for the prescribed edge. Indeed the top trace entry has $e_1=2(1+c_m)$, and the remaining even entries have $e_{k+1}=2(\Delta^k c)_{m-k}$ for $1\le k<m$.

Write $w_j=c_{m-j}\bmod2$. Modulo two, absolute differencing is addition of neighbors. Thus $b_j=\sum_{k=0}^j\binom{j}{k}w_k$ over $\mathbb F_2$. This triangular Pascal matrix is its own inverse because the $(j,k)$ entry of its square is $\binom jk 2^{j-k}$ for $k<j$ and is $1$ for $k=j$. Taking $j=m-1$ in the inversion proves necessity since $w_{m-1}=0$.

Conversely, reconstruct $c$ by the displayed inverse. Its first value is zero and every value is a bit. All entries of every difference row of $c$ are also bits, so the parities just computed give its exact integer trace. The positive gaps $2(1+c_i)$ define a strictly increasing prefix. The first row starts with $1$, and $c_1=0$ gives the second-row left entry $|2-1|=1$. In each later row the entry immediately to the right of the preceding $1$ is $2(\Delta^r c)_1\in\{0,2\}$ for the appropriate $r\ge1$. Its difference from $1$ is again $1$. This proves that the prefix is Gilbreath. Invertibility gives uniqueness within the asserted gap alphabet.
\end{proof}

\paragraph{The count at every length.}
Let $t=2^{\operatorname{popcount}(m-1)}$ and $z=m-t$, where $\operatorname{popcount}$ counts the ones in the binary expansion. There are $t$ odd binomial coefficients in row $m-1$, as follows by expanding $(1+X)^{m-1}$ over $\mathbb F_2$ into factors $1+X^{2^a}$. The number of distinct realizable right edges at widths $5,7,9$ is respectively
\[
 1,\qquad z,\qquad \binom{z}{2}+\binom{t}{2}.
\]
For weight one the selected index must have even coefficient. For weight two both selected indices must have the same coefficient parity. Width $7$ is absent exactly when $m$ is a power of two, including $m=1$. Indeed $t=m$ precisely when the binary digits of $m-1$ below its leading one are all ones, with the zero case treated separately. Width $9$ occurs for every $n\ge4$ since $t\ge2$ for $m\ge2$. At $n=3$ the unique prefix $(2,3,5)$ has width $5$. These counts enumerate right edges. Distinct prefixes with larger gaps can share an edge.


For example, at $n=7$ row $m-1=4$ of Pascal's triangle is odd at its two endpoints. Choosing $P=\{2\}$ gives $c=(0,1,0,1,0)$ and the prefix $(2,3,5,9,11,15,17)$ of width seven. At $n=10$ all eight entries of Pascal row seven are odd, so width seven is impossible. At $n=11$ it reappears, with seven distinct edges. Direct enumeration at lengths $3$ through $11$ gives the following comparison.
\[
\begin{array}{c|rrrrrrrrr}
 n&3&4&5&6&7&8&9&10&11\\\hline
 \text{distinct edges at width seven}&0&0&1&0&3&2&3&0&7\\
 \text{prefixes at width seven}&0&0&1&0&4&2&7&0&56
\end{array}
\]

\paragraph{Recovering every integer realization.}
For completeness, a finite branching construction also enumerates all such prefixes for any one of the permitted edges. Prescribe $u_{m-r,r}=b_r$ for $0\le r<m$. For $i=m-1,m-2,\ldots,1$, process $r=m-i-1,m-i-2,\ldots,0$, choosing each distinct nonnegative value
\[
 u_{i,r}\in\{u_{i+1,r}+u_{i,r+1},\ u_{i+1,r}-u_{i,r+1}\}.
\]
Keep precisely the arrays with $u_{1,0}=0$ and $u_{1,r}\in\{0,1\}$ for $1\le r<m$, and use $c_i=u_{i,0}$ in the gap formula. Each choice solves exactly $|u_{i+1,r}-u_{i,r}|=u_{i,r+1}$. The retained left-edge constraints are necessary and sufficient for the successive differences against $1$ in the original triangle to be $1$. Thus every surviving array gives one realization, and every realization has its own unique integer difference array.

\section{Quadratic logarithmic growth of the prefix count}
Let $\mathcal G_n$ denote the family of strictly increasing Gilbreath prefixes of length $n$ starting with $(2,3)$, and let $N_n=|\mathcal G_n|$. Translation by $-1$ identifies this count with OEIS A080839~\cite{noe}. In 2007 T. D. Noe recorded evidence for a limit near $1.46$ for $N_{n+1}N_{n-1}/N_n^2$. The following gives the scale of the logarithm without presupposing a limit.

\begin{theorem}
For all sufficiently large $n$,
\[
 2^{(\lfloor n/32\rfloor-1)\lfloor n/5\rfloor}
 \le N_n\le 2^{(n-1)(n-2)/2}.
\]
In particular $\log N_n=\Theta(n^2)$.
\end{theorem}

\begin{proof}
For the upper bound, start at the bottom entry $1$ and reconstruct the positive-depth rows upwards, always starting a row with $1$. Each further entry has at most two choices given its predecessor and their absolute difference. There are $(n-1)(n-2)/2$ such positions. Strict increase and the first entry $2$ fix the top row uniquely from its gaps.

For the lower bound set $h=\lfloor n/32\rfloor$ and prescribe row $h$ to be $(1,2,\ldots,2)$. Its descendants all have first entry $1$. Given each completed row, construct the row above it by the following rule, down to depth $1$. If the row below prescribes positive steps $a_1,\ldots,a_{L-1}$, set $y_1=1$. At position $i$ use only $y_i+a_i$ when $y_i<2a_i$ and allow both $y_i\pm a_i$ otherwise. The entries stay positive and the two choices are distinct whenever allowed. Finally start the top row at $2$ and add the successive entries of row $1$.

At offset $d$ rows above the prescribed row, the entry at position $i$ is bounded above by
\[
 b_d(i)=2\sum_{t=0}^{d}\binom{i-1}{t}.
\]
Indeed $b_0(i)=2$, $b_d(1)=2$, and $b_d(i+1)=b_d(i)+b_{d-1}(i)$ by Pascal's identity. In particular $y_L\le 2\sum_{t\le h}\binom{n}{t}\le 2^{1+n H_2(1/32)}$, where $H_2(x)=-x\log_2 x-(1-x)\log_2(1-x)$ is binary entropy. The binomial sum bound follows by summing the probabilities $\binom nt a^t(1-a)^{n-t}$ for $t\le an$, with $a=1/32$. Write $B$ for the number of positions allowing both signs along a completed row and let $q=\log_2(3/2)$. Each forced step increases the entry by a factor greater than $3/2$, and each free step by a factor at least $1/2$. Therefore
\[
 B\ge\frac{q(L-1)-\log_2 y_L}{1+q}
 \ge\frac{q(n-h)-1-nH_2(1/32)}{1+q}>\frac n5
\]
for sufficiently large $n$; here $L-1\ge n-h$ for every constructed row. In the binary choice tree for one row, give each free child half the probability of its parent and each forced child the full probability. Every leaf has mass at most $2^{-\lfloor n/5\rfloor}$, so the row admits at least $2^{\lfloor n/5\rfloor}$ possibilities. Multiply over the $h-1$ rows. Distinct triangles yield distinct top rows because the triangle is recovered by differencing its top row.
\end{proof}

\section{An arithmetic-progression conjecture of Kimberling}
\label{sec:kim}
For an odd positive integer $k$, let $S_k$ be the increasing sequence of all primes of the form $2+kt$ with $t\ge0$. It begins with $2$. In the comments to OEIS A358691, further conjectured example (1), Kimberling proposed that all positive-depth western entries of $S_k$ equal $k$ for every odd $k$ \cite{kim}. This universal assertion includes $k=1$ as just one case.

\begin{theorem}\label{thm:kim}
The proposed assertion is false at $k=7$. Moreover, writing $G_1(k)$ for its first western entry, one has
\[
 \frac{G_1(k)}{k}>\frac{\log k}{\log\log k}
\]
on a relative-density-one set of odd positive integers. For every $H>0$ there are also infinitely many odd $k$ for which $G_1(k)=k$ and the second western entry is greater than $Hk$.
\end{theorem}
\begin{proof}
At $k=7$ the candidates start with $2,9,16,23$. The first gap is $21$. In general put $M(k)=\min\{t\ge1:2+kt\text{ is prime}\}$, finite by Dirichlet's theorem, so $G_1(k)=kM(k)$. Only odd $t$ contribute. For $X\ge3$, $1\le L\le X$, a union bound and the Brun--Titchmarsh inequality \cite{bt} give
\[
 \#\{k\le X:k\text{ odd},\ M(k)\le L\}
 \le\sum_{\substack{t\le L\\t\ \mathrm{odd}}}\pi(tX+2;t,2)
 \ll \frac{X}{\log X}\sum_{t\le L}\frac{t}{\varphi(t)}
 \ll \frac{LX}{\log X}.
\]
Here $t=1$ is covered by the usual bound on $\pi(X+2)$. For the last inequality use $t/\varphi(t)=\sum_{d\mid t}\mu^2(d)/\varphi(d)$ and the convergent product $\prod_{\ell\ {\rm prime}}(1+1/[\ell(\ell-1)])$. Take $L=\lfloor\log X/\log\log X\rfloor$. The exceptional count is $o(X)$, and $\log k/\log\log k$ is eventually increasing. This gives the relative-density assertion.

Fix an integer $L>H+2$. For each odd $t$ with $3\le t\le L$, choose a distinct prime $\ell_t>L+2$. Prescribe an odd prime $q$ by the congruences
\[
 q\equiv 2(t-1)t^{-1}\pmod{\ell_t}.
\]
Every specified class is nonzero. The Chinese remainder theorem and Dirichlet's theorem therefore give infinitely many such $q$. For sufficiently large $q$ put $k=q-2$. The second prime in $S_k$ is $q$. Each $2+kt=tq-2(t-1)$ with odd $3\le t\le L$ is composite, being larger than its divisor $\ell_t$. The even indices also give composites. If the third prime in $S_k$ is $2+ku$, then $u>L$. The first two gaps are $k$ and $(u-1)k$, and the western entry at depth two is $(u-2)k>Hk$.
\end{proof}

The divisibility argument also describes what survives of the proposal at all moduli.
\begin{proposition}
For every odd $k\ge1$ the greatest common divisor of every western tail for $S_k$ is $k$. Consequently any eventual constant on that edge would have to equal $k$.
\end{proposition}
\begin{proof}
All positive-height entries are multiples of $k$. Normalize the top row to the integer indices $(p-2)/k$ for $p\in S_k$; its positive-height entries are the original ones divided by $k$. Fix a prime $\ell$. Choose $a\equiv2\pmod k$ coprime to $k\ell$ (at $\ell=2$ take $a$ odd). Shiu supplies arbitrarily long strings of consecutive primes all congruent to $a\pmod{k\ell}$; when $k\ell=2$, odd primes suffice. They are consecutive also in $S_k$, and their normalized indices all agree modulo $\ell$. If $\ell$ divided a whole tail of the normalized western edge, the same propagation and descent used in Section~\ref{sec:west} would imply divisibility of every normalized first gap. Its top row starts at zero, so all its indices would be multiples of $\ell$.

To contradict this, choose $t_0\not\equiv0\pmod\ell$ such that $2+kt_0\not\equiv0\pmod\ell$. For $\ell=2$ take $t_0=1$. For odd $\ell\mid k$ any nonzero $t_0$ suffices. For odd $\ell\nmid k$ at most two residues are forbidden. The class $2+kt_0$ is then coprime to $k\ell$ because $k$ is odd. Dirichlet gives arbitrarily large primes in that class modulo $k\ell$, whose normalized indices are nonzero modulo $\ell$. Hence no prime divides the normalized tail gcd.
\end{proof}

\section{What remains at the western edge}
The exact bound still needed for the prime sequence is $A_{r,2}\le2$ for all $r\ge1$. The tail gcd determines the value of a possible eventual constant, while the parity theorems force variation on interior columns and on moving antidiagonals. The increasing prime subsequence in Section~\ref{sec:west} preserves all interior residues modulo four and can depart from any prescribed prime prefix at the next step. This shows where an argument based on these invariants needs a further estimate specific to consecutive primes.

\smallskip
\noindent\textbf{Status of the individual questions.}
For the classical problem, the name, source and original quantifier are fixed by \eqref{eq:gil}. The universal magnitude estimate and hence a complete resolution remain open. For the finite problem in \cite{muney,vmwidth}, Theorem~\ref{thm:reachable} gives both directions and a realizing prefix at each permitted length. For Kimberling's further example (1), Theorem~\ref{thm:kim} decides the stated universal assertion by counterexample. The growing lower bound at the first step and the obstruction after a correct first step describe the extent of the failure. Example (2) in that entry is a separate eventual-constancy question.

\smallskip
\noindent\textbf{Attribution and priority.}
The Pascal involution, automatic-sequence criteria, prime-cluster theorems, and the natural-boundary dichotomy used here are established tools. The corpus already contains a classification by candidate antidiagonals and results on the maximum number of components. A separate public record also reports a solution of Muney's maximum-width Conjecture~30 and a proof of the component maximum \cite{components}. The literature check covered the cited prime-triangle papers, the OEIS entries, the relevant corpus sections and targeted searches for the ray, tail-gcd and realization results. This identifies the earlier questions and the mathematical steps added here. Historical priority of these proofs has not been established.

\begin{thebibliography}{99}\small
\bibitem{bftb} W. D. Banks, T. Freiberg and C. L. Turnage-Butterbaugh, Consecutive primes in tuples, \emph{Acta Arith.} \textbf{167} (2015), 261--266. \href{https://arxiv.org/abs/1311.7003}{arXiv:1311.7003}.
\bibitem{bcz} R. N. Bhat, C. Cobeli and A. Zaharescu, On quasi-periodicity in Proth--Gilbreath triangles, \emph{Bull. Math. Soc. Sci. Math. Roumanie} \textbf{67(115)} (2024), 3--21. \href{https://rms.unibuc.ro/bulletin/pdf/67-1/articol_1.pdf}{Conjecture 1 and Table 1}.
\bibitem{bcz2} R. N. Bhat, C. Cobeli and A. Zaharescu, Filtered rays over iterated absolute differences on layers of integers. \href{https://arxiv.org/abs/2309.03922}{arXiv:2309.03922}, Conjecture 2.
\bibitem{pc} F. Carlson, Über Potenzreihen mit ganzzahligen Koeffizienten, \emph{Math. Z.} \textbf{9} (1921), 1--13.
\bibitem{cht} Z. Chase, Z. Hunter and T. Tao, Gilbreath's conjecture: a Cramér random model and a deterministic analysis (2026). \href{https://arxiv.org/abs/2607.08712}{arXiv:2607.08712}.
\bibitem{christol} G. Christol, T. Kamae, M. Mendès France and G. Rauzy, Suites algébriques, automates et substitutions, \emph{Bull. Soc. Math. France} \textbf{108} (1980), 401--419. \href{https://doi.org/10.24033/bsmf.1926}{doi:10.24033/bsmf.1926}.
\bibitem{colonna} J.-F. Colonna, Proth--Gilbreath conjecture verification up to $1.5\times10^{15}$, computational report, 18 March 2026. \href{https://www.lactamme.polytechnique.fr/Mosaic/descripteurs/Proth_Gilbreath_Conjecture.01.Ang.html}{Report and verification range}.
\bibitem{kim} C. Kimberling, Gilbreath transform of odd-indexed primes, \href{https://oeis.org/A358691/internal}{OEIS A358691}, comments, further conjectured examples (1) and (2). Entry credited 27 November 2022.
\bibitem{noe} J. W. Layman, Number of increasing sequences with unit Gilbreath transform, \href{https://oeis.org/A080839}{OEIS A080839}. Growth observation by T. D. Noe, 5 February 2007.
\bibitem{maynard} J. Maynard, Dense clusters of primes in subsets, \emph{Compos. Math.} \textbf{152} (2016), 1517--1554. \href{https://doi.org/10.1112/S0010437X16007296}{doi:10.1112/S0010437X16007296}.
\bibitem{bt} H. L. Montgomery and R. C. Vaughan, The large sieve, \emph{Mathematika} \textbf{20} (1973), 119--134. \href{https://doi.org/10.1112/S0025579300004708}{doi:10.1112/S0025579300004708}.
\bibitem{muney} L. Muney, Holes in valid-extension sets of finite Gilbreath sequences, version 2, 2026, Sections 8--16. \href{https://arxiv.org/abs/2606.23721}{arXiv:2606.23721}.
\bibitem{odlyzko} A. M. Odlyzko, Iterated absolute values of differences of consecutive primes, \emph{Math. Comp.} \textbf{61} (1993), 373--380. \href{https://doi.org/10.1090/S0025-5718-1993-1182247-7}{doi:10.1090/S0025-5718-1993-1182247-7}.
\bibitem{proth} F. Proth, Sur la série des nombres premiers, \emph{Nouvelle Correspondance Mathématique} \textbf{4} (1878), 236--240. \href{https://gdz.sub.uni-goettingen.de/download/pdf/PPN598948236_0004/LOG_0088.pdf}{Original article}.
\bibitem{puchta} J.-C. Schlage-Puchta, A criterion for non-automaticity of sequences, \emph{J. Integer Seq.} \textbf{6} (2003), Article 03.3.8. \href{https://cs.uwaterloo.ca/journals/JIS/VOL6/Puchta/puchta70.pdf}{Article}.
\bibitem{shiu} D. K. L. Shiu, Strings of congruent primes, \emph{J. London Math. Soc.} (2) \textbf{61} (2000), 359--373. \href{https://doi.org/10.1112/S0024610799007863}{doi:10.1112/S0024610799007863}.
\bibitem{components} TrureTuring, Finite Gilbreath extension component maximum. \href{https://www.trureturing.org/research/muney-gilbreath-component-maximum/}{Public problem record}, which includes a pointer to S. Patel's maximum-cardinality result, MathDB 375513.
\bibitem{vmthree} Keio, \emph{Integrated Structures and Boundary Spectra in Sparse Prime Codings}, VMRC 003, section on fixed-depth Gilbreath obstruction languages. \href{https://github.com/QEEDDD/Verifiable-Mathematical-Reasoning-Corpus/tree/c2ac34ea4622778187173f41bb1996deee8081b7}{Repository revision c2ac34ea4622}.
\bibitem{vmwidth} Keio, \emph{Recurrences, Absolute Differences, and Collatz Orbits}, VMRC 048, also recorded under 072. Part B, subsection on widths at most nine, especially the remark on reachability. \href{https://github.com/QEEDDD/Verifiable-Mathematical-Reasoning-Corpus/tree/c2ac34ea4622778187173f41bb1996deee8081b7}{Repository revision c2ac34ea4622}.
\end{thebibliography}
\end{document}
